% ── chemistry-handout-preamble.tex ── 化学竞赛讲义 LaTeX 导言区（v7 · 精排版）
\documentclass[a4paper,11pt]{ctexart}

% ── 字体优化 ──
\setCJKmainfont[AutoFakeBold=true,BoldFont=SimHei]{SimSun}
\setCJKsansfont{SimHei}

% 化学式宏（mhchem）
\usepackage[version=4]{mhchem}

% ── 数学 ──
\usepackage{amsmath,amssymb}

% ── 图片 ──
\usepackage{graphicx}
\usepackage[font=small,labelfont=bf,skip=4pt]{caption}

% ── 表格（跨页+自动列宽）──
\usepackage{booktabs}
\usepackage{xltabular}
\usepackage{array}
\newcolumntype{L}{>{\raggedright\arraybackslash}X}
\newcolumntype{C}{>{\centering\arraybackslash}X}
\newcolumntype{R}{>{\raggedleft\arraybackslash}X}

% ── 浮动体（[H] 强制定位）──
\usepackage{float}

% ── 页面 ──
\usepackage[includeheadfoot,margin=2.2cm,headheight=14pt]{geometry}
\setlength{\textwidth}{16.6cm}

% ── 行间距 ──
\usepackage{setspace}
\setstretch{1.15}
\setlength{\parskip}{0.3em}
\setlength{\parindent}{2em}

% ── 防止表格跨页 ──
\usepackage{longtable}
\setlength{\LTpre}{0pt}
\setlength{\LTpost}{0pt}

% ── 防止图片跨页（尽量）──
\renewcommand{\textfraction}{0.1}
\renewcommand{\topfraction}{0.9}
\renewcommand{\bottomfraction}{0.9}
\setcounter{totalnumber}{10}
\setcounter{topnumber}{10}
\setcounter{bottomnumber}{10}
\setlength{\textfloatsep}{10pt plus 2pt minus 2pt}
\setlength{\floatsep}{8pt plus 2pt minus 2pt}
\setlength{\intextsep}{10pt plus 2pt minus 2pt}

% ── 超链接(无红框) ──
\usepackage[colorlinks=true,linkcolor=black,citecolor=black,urlcolor=black]{hyperref}

% ── 颜色 ──
\usepackage[dvipsnames,svgnames]{xcolor}
\definecolor{chapterblue}{RGB}{23,50,77}
\definecolor{insightbg}{RGB}{235,240,250}
\definecolor{warningbg}{RGB}{255,245,235}
\definecolor{tipbg}{RGB}{235,250,235}

% ── 页眉页脚 ──
\usepackage{fancyhdr}
\pagestyle{fancy}
\fancyhf{}
\fancyhead[L]{\small\color{chapterblue}\nouppercase{\leftmark}}
\fancyhead[R]{\small\color{chapterblue}\thepage}
\renewcommand{\headrule}{{\color{chapterblue}\hrule height 0.8pt}}
% 去掉页眉中的"超级充实版"等后缀
\renewcommand{\sectionmark}[1]{\markboth{#1}{}}
\patchcmd{\section}{\markright{\thesection\ --\ #1}}{\markright{\thesection}}{}{}

% ── 章节编号与样式 ──
\setcounter{secnumdepth}{1}
\ctexset{
  section = {
    format = \Large\bfseries\heiti\color{chapterblue},
    aftername = \hspace{0.5em},
    beforeskip = 2ex,
    afterskip = 1ex,
  },
  subsection = {
    format = \large\bfseries\heiti,
    beforeskip = 1.5ex,
    afterskip = 0.5ex,
  },
  subsubsection = {
    format = \normalsize\bfseries\heiti,
    beforeskip = 1ex,
    afterskip = 0.3ex,
  },
  paragraph = {
    format = \small\bfseries,
    beforeskip = 0.5ex,
    afterskip = 0.2ex,
  },
}

% ── blockquote（引用/提示框）缩减 ──
\usepackage{etoolbox}
\AtBeginEnvironment{quote}{\small\setstretch{1.05}\leftskip=1em\rightskip=1em}

% ── 教学提示框（浅底色版）──
\newenvironment{insightbox}{%
  \par\smallskip
  \noindent{\small\bfseries\color{RoyalBlue}🧠 认知冲突}\par
  \leavevmode\colorbox{insightbg}{\parbox{\dimexpr\textwidth-2\fboxsep\relax}{%
  \ignorespaces}}%
}{\par\smallskip}

\newenvironment{warningbox}{%
  \par\smallskip
  \noindent{\small\bfseries\color{BrickRed}⚠️ 易错点}\par
  \leavevmode\colorbox{warningbg}{\parbox{\dimexpr\textwidth-2\fboxsep\relax}{%
  \ignorespaces}}%
}{\par\smallskip}

\newenvironment{tipbox}{%
  \par\smallskip
  \noindent{\small\bfseries\color{ForestGreen}💡 理解提示}\par
  \leavevmode\colorbox{tipbg}{\parbox{\dimexpr\textwidth-2\fboxsep\relax}{%
  \ignorespaces}}%
}{\par\smallskip}

% ── 列表（答案区加间距）──
\usepackage{enumitem}
\setlist{nosep}
% 专门用于参考答案的列表（每项之间空一行）
\newlist{answerlist}{enumerate}{1}
\setlist[answerlist]{label=\arabic*., leftmargin=2em, itemsep=0.8em}

% ── 防止内容溢出 ──
\usepackage{microtype}
\setlength{\emergencystretch}{2em}
\setlength{\hfuzz}{2pt}
\setlength{\overfullrule}{0pt}
\sloppy

% ── pandoc 兼容 ──
\providecommand{\tightlist}{\setlength{\itemsep}{0pt}\setlength{\parskip}{0pt}}
\newenvironment{Shaded}{}{}
\newenvironment{Highlighting}{}{}
\newcommand{\NormalTok}[1]{#1}
\newcommand{\CommentTok}[1]{#1}
\newcommand{\KeywordTok}[1]{#1}
\newcommand{\StringTok}[1]{#1}
\newcommand{\FunctionTok}[1]{#1}
\newcommand{\DataTypeTok}[1]{#1}
\newcommand{\NumberTok}[1]{#1}
\newcommand{\ControlFlowTok}[1]{#1}
\newcommand{\OperatorTok}[1]{#1}
\newcommand{\BaseNTok}[1]{#1}
\newcounter{none}

% ── 图片尺寸限制 ──
\AtBeginDocument{\setkeys{Gin}{width=\textwidth,keepaspectratio}}

% ── 防止图片跨节 ──
\usepackage{placeins}

% ── 自定义命令 ──
% 练习题标题（去掉括号）
% 例题标题（去掉中括号）
% 这些需要在convert.py预处理中做正则替换

% ── Unicode 符号渲染命令（preprint 中调用，\ensuremath 避免 $...$ 切断 CJK 字间距）──
\newcommand{\箭}{\ensuremath{\rightarrow}}
\newcommand{\左箭}{\ensuremath{\leftarrow}}
\newcommand{\上箭}{\ensuremath{\uparrow}}
\newcommand{\下箭}{\ensuremath{\downarrow}}
\newcommand{\衡}{\ensuremath{\rightleftharpoons}}
\newcommand{\双箭}{\ensuremath{\leftrightarrow}}
\newcommand{\约等于}{\ensuremath{\approx}}
\newcommand{\小于等于}{\ensuremath{\leq}}
\newcommand{\大于等于}{\ensuremath{\geq}}
\newcommand{\不等于}{\ensuremath{\neq}}
\newcommand{\西格马}{\ensuremath{\sigma}}
\newcommand{\派}{\ensuremath{\pi}}

\graphicspath{{C:/Obsidion/妙妙屋/11-模板/scripts/.chem_media/}}
\title{元素周期表与周期律}
\date{}
\setlength{\tabcolsep}{3.5pt}
\small
\begin{document}
\thispagestyle{plain}
\begin{center}
{\zihao{2}\bfseries 元素周期表与周期律\par}
\end{center}
\vspace{1em}
\section{元素周期表与元素周期律：超级充实版}\label{ux5143ux7d20ux5468ux671fux8868ux4e0eux5143ux7d20ux5468ux671fux5f8bux8d85ux7ea7ux5145ux5b9eux7248}

\subsection{学习目标}\label{ux5b66ux4e60ux76eeux6807}

\begin{itemize}
\tightlist
\item
  能根据鲍林能级组划分解释元素周期表的周期结构，判断元素所在的周期、族和区
\item
  能比较原子半径大小(三种半径定义的选择与判断)
\item
  能运用有效核电荷和电子排布原理系统解释原子半径、电离能、电子亲和能、电负性的周期性变化趋势
\item
  能准确指出并解释周期律中的关键反常点(IIA/IIIA、VA/VIA电离能倒置；F/Cl电子亲和能反常)
\item
  能用镧系收缩解释第二、第三过渡系同族元素性质相似的原因
\item
  能用对角线规则解释 Li-Mg、Be-Al、B-Si 的化学相似性
\item
  能通过电离能、电负性等数据推断未知元素在周期表中的位置与族属
\item
  能用 Allred-Rochow 电负性标度定性理解 Z* 与电负性的关系
\item
  能用惰性电子对效应解释 Tl(I)/Pb(II)/Bi(III) 的稳定性
\item
  能用交换能定性解释半满/全满稳定化及 Cr/Cu 等异常排布
\item
  能用次级周期性解释第四/六周期高价化合物氧化性增强
\item
  理解相对论效应对重元素(Au 颜色、Hg 液态、惰性电子对)的统一解释
\end{itemize}

\subsection{一、元素周期表的结构}\label{ux4e00ux5143ux7d20ux5468ux671fux8868ux7684ux7ed3ux6784}

\begin{quote}
\textbf{理解要点}：在上一讲我们学会了用四个量子数给轨道贴标签、用三原则给电子排布。但排布本身不是终点---\textbf{排布决定了元素的性质}。周期表是''排布的目录''，周期律是''性质变化的规律''。有了这个工具，你才能：
- 看到一种元素，立刻判断它的金属性/非金属性、常见的氧化态、主要键型
- 看到一组电离能数据，反推出元素在周期表中的位置
- 理解为什么Zr和Hf分离困难、为什么Pb(II)比Pb(IV)稳定、为什么Li和Mg性质相似

一句话：\textbf{原子结构回答''电子怎么排''，周期律回答''排好之后怎么样''}，两者加起来才是你判断一切元素性质的基础。
\end{quote}

\subsubsection{1.1 周期---能级组的划分}\label{ux5468ux671fux80fdux7ea7ux7ec4ux7684ux5212ux5206}

元素周期表的周期划分依据是鲍林近似能级图中的\textbf{能级组}。每个能级组对应一个周期，能级组内轨道数×2 = 该周期元素数目。

\textbf{周期速查}：第一～三周期为短周期(2/8/8 种)，第四～五为长周期(18 种)，第六～七为特长周期(32 种，含 f 区)。周期数 = 最外层电子层主量子数 \(n\)。

\textbf{元素数目规律}：第四能级组含 \(4\mathrm{s}(1)+3\mathrm{d}(5)+4\mathrm{p}(3)=9\) 个轨道 \箭 \(9\times2=18\) 种元素。以此类推。

\begin{figure}[H]\centering
\includegraphics[width=0.58\textwidth,keepaspectratio]{11-20-cotton-orbital-energy-vs-Z.jpg}
\caption{Cotton 轨道能量随原子序数 Z 的变化图。}
\end{figure}

\begin{quote}
\textbf{读图提示}：横轴为原子序数 Z，纵轴为轨道能量。注意 4s 与 3d 在 Z\约等于20 附近交叉(K/Ca 处 4s 低于 3d，Sc 之后 3d 迅速下降)，这就是为什么第四周期先填 4s 再填 3d。放大的右上角显示了 K\箭V 的填充顺序。
\end{quote}

\subsubsection{1.2 族---价电子数决定族属}\label{ux65cfux4ef7ux7535ux5b50ux6570ux51b3ux5b9aux65cfux5c5e}

周期表 18 列，划分为 \textbf{16 个族}：7 个主族(A)、7 个副族(B)、1 个零族(稀有气体)、1 个第VIII族(8、9、10 列)。

\textbf{主族元素}(s 区 + p 区)：族数 = 最外层价电子数(\(ns + np\) 电子数之和)。和为 8 时为零族(稀有气体)。

\textbf{副族元素}(d 区 + ds 区)：
- IIIB～VIIB 族：族数 = 价电子数(\((n-1)\mathrm{d} + n\mathrm{s}\) 之和)
- IB 族：价层电子构型 \((n-1)\mathrm{d}^{10}n\mathrm{s}^{1}\)
- IIB 族：价层电子构型 \((n-1)\mathrm{d}^{10}n\mathrm{s}^{2}\)
- 第VIII族(8、9、10 列)：价电子数 \大于等于 8

\textbf{过渡元素}：周期表中副族元素(IIIB～IIB)和第VIII族元素的总称。其共同特征是价电子排布中包含 \((n-1)\mathrm{d}\) 轨道。

\textbf{镧系与锕系}：第六周期IIIB 族位置上填有 \(_{57}\mathrm{La}\sim _{71}\mathrm{Lu}\) 共 15 种元素(镧系)；第七周期对应位置有 \(_{89}\mathrm{Ac}\sim _{103}\mathrm{Lr}\) 共 15 种元素(锕系)。镧系元素和 Sc、Y 统称为\textbf{稀土元素}。

\subsubsection{1.3 元素分区---价电子构型决定分区}\label{ux5143ux7d20ux5206ux533aux4ef7ux7535ux5b50ux6784ux578bux51b3ux5b9aux5206ux533a}

根据最后一个电子填充的能级类型，周期表元素分为五个区：

{\def\LTcaptype{none} % do not increment counter
\begin{longtable}[]{@{}
  >{\centering\arraybackslash}p{(\linewidth - 6\tabcolsep) * \real{0.2941}}
  >{\raggedright\arraybackslash}p{(\linewidth - 6\tabcolsep) * \real{0.2353}}
  >{\raggedright\arraybackslash}p{(\linewidth - 6\tabcolsep) * \real{0.2353}}
  >{\raggedright\arraybackslash}p{(\linewidth - 6\tabcolsep) * \real{0.2353}}@{}}
\toprule\noalign{}
\begin{minipage}[b]{\linewidth}\centering
区
\end{minipage} & \begin{minipage}[b]{\linewidth}\raggedright
价电子构型
\end{minipage} & \begin{minipage}[b]{\linewidth}\raggedright
包含族
\end{minipage} & \begin{minipage}[b]{\linewidth}\raggedright
特征
\end{minipage} \\
\midrule\noalign{}
\endhead
\bottomrule\noalign{}
\endlastfoot
\textbf{s 区} & \(n\mathrm{s}^{1\sim2}\) & IA、IIA & 最外层为 \(s\) 电子；碱金属和碱土金属 \\
\textbf{p 区} & \(n\mathrm{s}^{2}n\mathrm{p}^{1\sim6}\) & IIIA～VIIA、零族 & 最外层有 \(p\) 电子 \\
\textbf{d 区} & \((n-1)\mathrm{d}^{1\sim10}n\mathrm{s}^{0\sim2}\) & IIIB～VIIB、VIII族 & 过渡元素；最后填 \(d\) 电子 \\
\textbf{ds 区} & \((n-1)\mathrm{d}^{10}n\mathrm{s}^{1\sim2}\) & IB、IIB & \(d\) 轨道全满后填 \(s\) 电子 \\
\textbf{f 区} & \((n-2)\mathrm{f}^{0\sim14}(n-1)\mathrm{d}^{0\sim2}n\mathrm{s}^{2}\) & 镧系、锕系 & 内过渡元素；最后填 \(f\) 电子 \\
\end{longtable}
}

\textbf{金属/非金属/半金属分区}：p 区中存在一条从 B 到 At 的阶梯分界线，将金属与非金属分开。分界线附近的元素(B、Si、Ge、As、Sb、Te、At)称为\textbf{半金属}(metalloid)，兼具金属和非金属特征。

{\def\LTcaptype{none} % do not increment counter
\begin{longtable}[]{@{}lll@{}}
\toprule\noalign{}
分类 & 电负性范围 & 典型代表 \\
\midrule\noalign{}
\endhead
\bottomrule\noalign{}
\endlastfoot
\textbf{金属} & \(\chi < 2.0\) & Li、Na、Fe、Cu、Zn 等 \\
\textbf{半金属} & \(\chi \approx 1.8\text{-}2.2\) & B、Si、Ge、As、Sb、Te \\
\textbf{非金属} & \(\chi > 2.0\) & C、N、O、F、Cl、S 等 \\
\end{longtable}
}

\begin{quote}
\textbf{关联知识点}：元素周期律\#价电子构型与族的关系 · 原子结构-超级充实版(自学完整)\#3.2 能级组与周期表的对应关系
\end{quote}

\begin{quote}
\textbf{名师点拨}：区分 d 区和 ds 区的关键判据：只看最后一个电子填在哪个轨道是不够的，还要看 d 轨道是否全满。d 区 = 未满/半满，ds 区 = 全满 + 再加 s 电子。

快速记忆：ds = d + s，d 轨道先''住满''10个人，s 轨道再住1-2个。
\end{quote}

\begin{quote}
\textbf{判定方法}：写出基态原子电子排布 \箭 找出最后一个电子填入的能级 \箭 确定分区。示例：
- Fe (\([\mathrm{Ar}] 3\mathrm{d}^6 4\mathrm{s}^2\))：最后一个电子填入 \(3\mathrm{d}\) \箭 d 区 \箭 第四周期 VIII族
- Cu (\([\mathrm{Ar}] 3\mathrm{d}^{10} 4\mathrm{s}^1\))：最后一个电子填入 \(4\mathrm{s}\)，但 \(3\mathrm{d}\) 已全满 \箭 ds 区 \箭 第四周期 IB 族
\end{quote}

\begin{figure}[H]\centering
\includegraphics[width=0.55\textwidth,keepaspectratio]{periodic-table-zones.jpg}
\caption{周期表分区与价电子构型的对应关系。}
\end{figure}

\begin{quote}
\textbf{读图提示}：s区(第1-2列)、p区(第13-18列)、d区(第3-12列)、ds区(第11-12列)、f区(镧系+锕系)。注意 ds 区是 d 区的最右两列，d 轨道已全满。
\end{quote}

\begin{figure}[H]\centering
\includegraphics[width=0.55\textwidth,keepaspectratio]{diagonal-rule-bw.png}
\caption{对角线规则与电荷密度对比。}
\end{figure}

\begin{center}\rule{0.5\linewidth}{0.5pt}\end{center}

\subsection{二、元素周期律}\label{ux4e8cux5143ux7d20ux5468ux671fux5f8b}

\begin{quote}
\textbf{理解要点}：原子半径、电离能、电子亲和能、电负性是描述元素性质的四个最基本工具。把它们搞清楚了，你就可以：
- 看到任意两个元素，比较谁更容易失电子、谁更容易得电子
- 从数据中反推元素在周期表中哪里
- 理解为什么有的元素像金属、有的像非金属，有的''模棱两可''

本质就一句话：\textbf{四大参数 = 电子排布的外在表现}。
\end{quote}

元素的性质随原子序数递增而呈现周期性变化的规律称为\textbf{元素周期律}。其本质是核外电子排布的周期性。

\begin{quote}
\textbf{关联知识点}：元素周期律 · 原子半径 · 电离能 · 电负性
\end{quote}

\subsubsection{2.1 原子半径}\label{ux539fux5b50ux534aux5f84}

\textbf{三种半径定义}

{\def\LTcaptype{none} % do not increment counter
\begin{longtable}[]{@{}
  >{\raggedright\arraybackslash}p{(\linewidth - 4\tabcolsep) * \real{0.3333}}
  >{\raggedright\arraybackslash}p{(\linewidth - 4\tabcolsep) * \real{0.3333}}
  >{\raggedright\arraybackslash}p{(\linewidth - 4\tabcolsep) * \real{0.3333}}@{}}
\toprule\noalign{}
\begin{minipage}[b]{\linewidth}\raggedright
类型
\end{minipage} & \begin{minipage}[b]{\linewidth}\raggedright
定义
\end{minipage} & \begin{minipage}[b]{\linewidth}\raggedright
适用对象
\end{minipage} \\
\midrule\noalign{}
\endhead
\bottomrule\noalign{}
\endlastfoot
\textbf{共价半径} & 同种元素以共价单键结合时，两原子核间距的一半 & 非金属元素 \\
\textbf{金属半径} & 金属晶体中相邻原子核间距的一半 & 金属元素 \\
\textbf{范德华半径} & 相邻分子间靠分子间作用力接触时，两相同原子核间距的一半 & 稀有气体 \\
\end{longtable}
}

同一元素三种半径的大小关系：\textbf{范德华半径 \textgreater{} 金属半径 \textgreater{} 共价半径}

\textbf{周期性变化规律}

\textbf{同周期}(从左到右)：\textbf{原子半径递减}(稀有气体除外)

本质：\(Z^*\) 增大 \箭 核对电子引力增强 \箭 半径减小。主族相邻平均减小约 10 pm，过渡元素仅约 5 pm(d 电子屏蔽较好)，镧系不足 1 pm(f 电子屏蔽更好)。

\textbf{同主族}(从上到下)：\textbf{原子半径递增}

本质：电子层数增加起主导作用，\(Z^*\) 增幅不明显。

\begin{figure}[H]\centering
\includegraphics[width=0.55\textwidth,keepaspectratio]{effective_nuclear_charge.png}
\caption{有效核电荷 Z* 示意图—Cl(Z*=+7，半径小)vs Mg(Z*=+2，半径大)。}
\end{figure}

\begin{quote}
\textbf{读图提示}：两原子电子层数相同(3 层)，但 Cl 的 17 个质子对外层电子的净吸引力(Z\emph{=+7)远强于 Mg 的 12 个质子(Z}=+2)，因此 Cl 半径更小。这就是同周期从左到右半径递减的物理图像。
\end{quote}

\textbf{1. 算一算：同周期半径递减的定量感受}

已知 Na 的原子半径 186 pm，Cl 的原子半径 99 pm。从 Na 到 Cl，核电荷从 11 增加到 17，半径却几乎减半。用 Slater 规则估算有效核电荷的变化：

\begin{itemize}
\tightlist
\item
  用 Slater 规则粗略估算 Na 的 \(Z^*\)(3s 电子)：\(\sigma \approx 1.00\times 2 + 0.85\times 8 = 8.8\)，\(Z^* \approx 11 - 8.8 = 2.2\)
\item
  估算 Cl 的 \(Z^*\)(3p 电子)：\(\sigma \approx 1.00\times 2 + 0.85\times 8 + 0.35\times 6 = 10.9\)，\(Z^* \approx 17 - 10.9 = 6.1\)
\item
  \(Z^*\) 从 2.2 增大到 6.1(约 2.8 倍)，半径从 186 pm 减小到 99 pm(约 0.53 倍)\箭 半径大致与 \(1/Z^*\) 成正比
\end{itemize}

这个估算让你直观感受到：\textbf{半径变化的主要驱动力不是核电荷多了几个，而是有效核电荷翻了几倍}。

\textbf{镧系收缩}

\textbf{定义}：镧系元素(\(_{57}\mathrm{La}\sim_{71}\mathrm{Lu}\))的原子半径和离子半径随原子序数增大而逐渐减小的现象。

\textbf{原因}：新增电子填入 \(4\mathrm{f}\) 轨道(无径向节点，高度紧缩)，\(f\) 电子对核电荷的屏蔽不充分 \箭 有效核电荷逐渐增大 \箭 核对最外层电子吸引力增强 \箭 半径逐渐减小。

\textbf{Ln\textsuperscript{3+} 离子半径数据}(pm)：

{\def\LTcaptype{none} % do not increment counter
\begin{longtable}[]{@{}
  >{\centering\arraybackslash}p{(\linewidth - 28\tabcolsep) * \real{0.0667}}
  >{\centering\arraybackslash}p{(\linewidth - 28\tabcolsep) * \real{0.0667}}
  >{\centering\arraybackslash}p{(\linewidth - 28\tabcolsep) * \real{0.0667}}
  >{\centering\arraybackslash}p{(\linewidth - 28\tabcolsep) * \real{0.0667}}
  >{\centering\arraybackslash}p{(\linewidth - 28\tabcolsep) * \real{0.0667}}
  >{\centering\arraybackslash}p{(\linewidth - 28\tabcolsep) * \real{0.0667}}
  >{\centering\arraybackslash}p{(\linewidth - 28\tabcolsep) * \real{0.0667}}
  >{\centering\arraybackslash}p{(\linewidth - 28\tabcolsep) * \real{0.0667}}
  >{\centering\arraybackslash}p{(\linewidth - 28\tabcolsep) * \real{0.0667}}
  >{\centering\arraybackslash}p{(\linewidth - 28\tabcolsep) * \real{0.0667}}
  >{\centering\arraybackslash}p{(\linewidth - 28\tabcolsep) * \real{0.0667}}
  >{\centering\arraybackslash}p{(\linewidth - 28\tabcolsep) * \real{0.0667}}
  >{\centering\arraybackslash}p{(\linewidth - 28\tabcolsep) * \real{0.0667}}
  >{\centering\arraybackslash}p{(\linewidth - 28\tabcolsep) * \real{0.0667}}
  >{\centering\arraybackslash}p{(\linewidth - 28\tabcolsep) * \real{0.0667}}@{}}
\toprule\noalign{}
\begin{minipage}[b]{\linewidth}\centering
La
\end{minipage} & \begin{minipage}[b]{\linewidth}\centering
Ce
\end{minipage} & \begin{minipage}[b]{\linewidth}\centering
Pr
\end{minipage} & \begin{minipage}[b]{\linewidth}\centering
Nd
\end{minipage} & \begin{minipage}[b]{\linewidth}\centering
Pm
\end{minipage} & \begin{minipage}[b]{\linewidth}\centering
Sm
\end{minipage} & \begin{minipage}[b]{\linewidth}\centering
Eu
\end{minipage} & \begin{minipage}[b]{\linewidth}\centering
Gd
\end{minipage} & \begin{minipage}[b]{\linewidth}\centering
Tb
\end{minipage} & \begin{minipage}[b]{\linewidth}\centering
Dy
\end{minipage} & \begin{minipage}[b]{\linewidth}\centering
Ho
\end{minipage} & \begin{minipage}[b]{\linewidth}\centering
Er
\end{minipage} & \begin{minipage}[b]{\linewidth}\centering
Tm
\end{minipage} & \begin{minipage}[b]{\linewidth}\centering
Yb
\end{minipage} & \begin{minipage}[b]{\linewidth}\centering
Lu
\end{minipage} \\
\midrule\noalign{}
\endhead
\bottomrule\noalign{}
\endlastfoot
116 & 114 & 113 & 111 & 109 & 108 & 107 & 106 & 104 & 103 & 102 & 100 & 99 & 99 & 98 \\
\end{longtable}
}

总收缩约 \textbf{18 pm}，相邻元素仅差约 \textbf{1 pm}。Ln\textsuperscript{3+} 离子半径严格单调下降(无例外)，因为所有 Ln\textsuperscript{3+} 都是 {[}Xe{]}4fⁿ 构型。

\textbf{Eu 和 Yb 金属半径例外}：Eu(4f\textsuperscript{7}6s\textsuperscript{2}) 和 Yb(4f\textsuperscript{14}6s\textsuperscript{2}) 的半满/全满构型特别稳定，金属键形成时不愿破坏此构型 \箭 金属键较弱 \箭 金属半径异常偏大(Eu 205 pm, Yb 195 pm)。但它们的 Ln\textsuperscript{3+} 离子半径仍遵循正常趋势。

\textbf{镧系收缩 vs d 区收缩}：

{\def\LTcaptype{none} % do not increment counter
\begin{longtable}[]{@{}lll@{}}
\toprule\noalign{}
对比项 & 镧系收缩 & d 区收缩 \\
\midrule\noalign{}
\endhead
\bottomrule\noalign{}
\endlastfoot
涉及元素 & 15 种(La\箭Lu) & 10 种(Sc\箭Zn 族) \\
总收缩量 & \textasciitilde18 pm & \textasciitilde5 pm \\
原因 & 4f 无径向节点，屏蔽极差 & 3d 有一定径向分布，屏蔽较好 \\
\end{longtable}
}

\textbf{核心意义}---镧系收缩是以下事实的根源：
1. \textbf{第二、第三过渡系同族元素半径相近}：Zr(160 pm)/Hf(159 pm)、Nb(146 pm)/Ta(146 pm)、Mo(139 pm)/W(139 pm)---半径几乎相等，化学性质极为相似，分离困难
2. \textbf{离子交换分离稀土}：半径越小的 Ln\textsuperscript{3+} 与配体形成的配合物越稳定 \箭 在离子交换柱上先被洗脱(反直觉！半径小=先出来，是竞赛常考陷阱)
3. \textbf{第三过渡系电离能普遍偏高}：半径与第二过渡系相近但核电荷大增

\subsubsection{2.2 电离能}\label{ux7535ux79bbux80fd}

\begin{quote}
\textbf{理解要点}：电离能告诉你两件事：一是这元素失电子有多容易(金属性强弱)，二是它有几个价电子(\箭判断族属)。它是元素推断题中最有力的数据工具之一。
\end{quote}

\textbf{定义}：基态气态原子失去一个电子成为一价气态正离子所需的最低能量，称为第一电离能 \(I_1\)。依次有 \(I_2, I_3, \dots\)。同一元素各级电离能：\(I_1 < I_2 < I_3 < \cdots\)(失去电子后 \(Z^*\) 增大，剩余电子受核吸引更强)。
\textbf{周期趋势}

\textbf{同周期}(从左到右)：\(I_1\) \textbf{总体递增}，稀有气体最大。

\textbf{同主族}(从上到下)：\(I_1\) \textbf{总体递减}(电子层数增加 \textgreater{} 核电荷增加)。

\textbf{副族}：第三过渡系 \(I_1\) 明显大于第二过渡系(镧系收缩 \箭 半径相近但核电荷大增)。

\textbf{关键反常点}

{\def\LTcaptype{none} % do not increment counter
\begin{longtable}[]{@{}
  >{\raggedright\arraybackslash}p{(\linewidth - 4\tabcolsep) * \real{0.3333}}
  >{\raggedright\arraybackslash}p{(\linewidth - 4\tabcolsep) * \real{0.3333}}
  >{\raggedright\arraybackslash}p{(\linewidth - 4\tabcolsep) * \real{0.3333}}@{}}
\toprule\noalign{}
\begin{minipage}[b]{\linewidth}\raggedright
反常
\end{minipage} & \begin{minipage}[b]{\linewidth}\raggedright
数据对比 (kJ/mol)
\end{minipage} & \begin{minipage}[b]{\linewidth}\raggedright
原因
\end{minipage} \\
\midrule\noalign{}
\endhead
\bottomrule\noalign{}
\endlastfoot
\textbf{IIA \textgreater{} IIIA} & Be(900) \textgreater{} B(801); Mg(738) \textgreater{} Al(578) & IIA 族 \(n\mathrm{s}^2\) 全满稳定；IIIA 族失一个 \(n\mathrm{p}\) 电子后得 \(n\mathrm{s}^2\) 稳定构型 \\
\textbf{VA \textgreater{} VIA} & N(1402) \textgreater{} O(1314); P(1012) \textgreater{} S(1000) & VA 族 \(n\mathrm{p}^3\) 半满稳定；VIA 族失一个电子后得 \(n\mathrm{p}^3\) 半满构型 \\
\textbf{Mn \textgreater{} 相邻元素} & Mn(717) \textgreater{} Cr(653) & Mn 的 \(3\mathrm{d}^5\) 半满构型稳定 \\
\textbf{Zn \textgreater{} 相邻元素} & Zn(906) \textgreater{} Cu(746) & Zn 的 \(3\mathrm{d}^{10}\) 全满构型稳定 \\
\end{longtable}
}

\begin{quote}
\textbf{点拨}：``比较电离能，先看轨道类型，再看周期趋势''。学生最常见的错误是只背''从左到右增大''，忽略了IIA/IIIA和VA/VIA两处关键反常。根源在于：\textbf{电子排布的稳定性对电离能的影响，有时大于核电荷增加的影响}。

\textbf{辨析}：``不要把周期律背成几个箭头---电离能的走势是'上台阶偶尔绊一下'，不是'上坡路'。''
\end{quote}

\textbf{交换能与半满/全满稳定}

为什么 \(n\mathrm{s}^2\) 全满和 \(n\mathrm{p}^3\) 半满格外稳定？根源在于\textbf{交换能}(Exchange Energy)。

\textbf{定性公式}：
\[E_{\text{ex}} \approx -K \cdot \sum \frac{n_i(n_i-1)}{2}\]

其中 \(n_i\) 为第 \(i\) 个亚层中\textbf{平行自旋}电子数，\(K\) 为交换积分(正值)。平行自旋电子对越多 \箭 交换能越负(体系越稳定)。

\textbf{实例}：Cr 和 Cu 的异常排布

{\def\LTcaptype{none} % do not increment counter
\begin{longtable}[]{@{}
  >{\centering\arraybackslash}p{(\linewidth - 6\tabcolsep) * \real{0.2500}}
  >{\centering\arraybackslash}p{(\linewidth - 6\tabcolsep) * \real{0.2500}}
  >{\centering\arraybackslash}p{(\linewidth - 6\tabcolsep) * \real{0.2500}}
  >{\centering\arraybackslash}p{(\linewidth - 6\tabcolsep) * \real{0.2500}}@{}}
\toprule\noalign{}
\begin{minipage}[b]{\linewidth}\centering
构型
\end{minipage} & \begin{minipage}[b]{\linewidth}\centering
平行自旋对数
\end{minipage} & \begin{minipage}[b]{\linewidth}\centering
交换能相对大小
\end{minipage} & \begin{minipage}[b]{\linewidth}\centering
实际选择
\end{minipage} \\
\midrule\noalign{}
\endhead
\bottomrule\noalign{}
\endlastfoot
\(3\mathrm{d}^4 4\mathrm{s}^2\) & \(\binom{4}{2}=6\) & 较小 & ✗ \\
\(3\mathrm{d}^5 4\mathrm{s}^1\) & \(\binom{5}{2}=10\) & \textbf{较大} & ✓ Cr 选此构型 \\
\(3\mathrm{d}^9 4\mathrm{s}^2\) & \(\binom{5}{2}+\binom{4}{2}=16\) & 中等 & ✗ \\
\(3\mathrm{d}^{10} 4\mathrm{s}^1\) & \(\binom{5}{2}\times2=20\) & \textbf{最大} & ✓ Cu 选此构型 \\
\end{longtable}
}

\begin{quote}
\textbf{要点}：交换能收益 + 对称性稳定化 \textgreater{} 能级差时，就会发生''异常''排布。这个逻辑同样解释了为什么 VA 族的 \(n\mathrm{p}^3\) 半满比 VIA 族的 \(n\mathrm{p}^4\) 更稳定---多了 \(\binom{3}{2}=3\) 个平行自旋对的交换能收益，同时少了配对电子的排斥能。

\textbf{关联知识点}：交换能的完整讨论见 原子结构-超级充实版(自学完整)\#2.2 交换能。
\end{quote}

\textbf{电离能的半定量估算}

利用类氢原子模型可以半定量估算第一电离能：
\[I_1 \approx 13.6 \cdot \frac{(Z^*)^2}{n^2} \quad (\text{eV})\]

例如：Na(\(Z^*\approx2.2\)，\(n=3\))：\(I_1 \approx 13.6 \times 2.2^2/9 \approx 7.3\) eV \约等于 704 kJ/mol(实测 496 kJ/mol)。虽然绝对值偏差较大(因为是粗略模型)，但\textbf{同周期趋势和反常方向}能正确再现---半定量公式的真正价值在于建立物理直觉，而非精确计算。

\textbf{1. 算一算：用电离能数据推断族属}

已知某元素 \(I_1\sim I_5\) (kJ/mol) = 786, 1580, 3230, 4360, 16000

\textbf{四步推断法}：
1. 计算各步增幅：\(I_2/I_1=2.0\times\)，\(I_3/I_2=2.0\times\)，\(I_4/I_3=1.4\times\)，\(I_5/I_4=3.7\times\)
2. 找到剧增点：\(I_4\to I_5\) 突增 \箭 失去 4 个电子后达到稳定构型
3. 价电子数 = 剧增前已失去的电子数 = 4 \箭 **IVA 族**
4. 若该元素为第三周期 \箭 Si(\([\mathrm{Ne}]3\mathrm{s}^2 3\mathrm{p}^2\))

\textbf{练习}：某元素 \(I_1\sim I_4\) = 578, 1817, 2745, 11575，推断族属。(答案见§七)

\textbf{电离能突变的竞赛应用}

电离能数据可作为元素推断的核心工具：\textbf{\(I_n \to I_{n+1}\) 剧增处 \箭 失去第 \(n\) 个电子后达到稀有气体稳定构型 \箭 该元素有 \(n\) 个价电子 \箭 族属可确定}。

示例：某元素 \(I_1\sim I_4\) (kJ/mol) = 580, 1815, 2740, 11600
- \(I_3\to I_4\) 剧增(4.2 倍)\箭 失去 3 个电子后达稳定构型 \箭 3 个价电子 \箭 **IIIA 族**

\textbf{典型元素逐级电离能(kJ/mol)}

掌握几个典型元素的逐级 IE 数据，有助于在竞赛中快速判断元素类型：

{\def\LTcaptype{none} % do not increment counter
\begin{longtable}[]{@{}
  >{\centering\arraybackslash}p{(\linewidth - 14\tabcolsep) * \real{0.1250}}
  >{\centering\arraybackslash}p{(\linewidth - 14\tabcolsep) * \real{0.1250}}
  >{\centering\arraybackslash}p{(\linewidth - 14\tabcolsep) * \real{0.1250}}
  >{\centering\arraybackslash}p{(\linewidth - 14\tabcolsep) * \real{0.1250}}
  >{\centering\arraybackslash}p{(\linewidth - 14\tabcolsep) * \real{0.1250}}
  >{\centering\arraybackslash}p{(\linewidth - 14\tabcolsep) * \real{0.1250}}
  >{\centering\arraybackslash}p{(\linewidth - 14\tabcolsep) * \real{0.1250}}
  >{\centering\arraybackslash}p{(\linewidth - 14\tabcolsep) * \real{0.1250}}@{}}
\toprule\noalign{}
\begin{minipage}[b]{\linewidth}\centering
元素
\end{minipage} & \begin{minipage}[b]{\linewidth}\centering
\(I_1\)
\end{minipage} & \begin{minipage}[b]{\linewidth}\centering
\(I_2\)
\end{minipage} & \begin{minipage}[b]{\linewidth}\centering
\(I_3\)
\end{minipage} & \begin{minipage}[b]{\linewidth}\centering
\(I_4\)
\end{minipage} & \begin{minipage}[b]{\linewidth}\centering
\(I_5\)
\end{minipage} & \begin{minipage}[b]{\linewidth}\centering
价电子数
\end{minipage} & \begin{minipage}[b]{\linewidth}\centering
突增点
\end{minipage} \\
\midrule\noalign{}
\endhead
\bottomrule\noalign{}
\endlastfoot
Na & 496 & 4562 & 6910 & 9543 & 13354 & 1 & \(I_1\to I_2\)(9.2×) \\
Mg & 738 & 1451 & 7733 & 10540 & 13630 & 2 & \(I_2\to I_3\)(5.3×) \\
Al & 578 & 1817 & 2745 & 11577 & 14842 & 3 & \(I_3\to I_4\)(4.2×) \\
Fe & 762 & 1561 & 2957 & 5290 & 7240 & 2-3 & \(I_2\to I_3\) 无剧增(3d 电子) \\
\end{longtable}
}

\begin{quote}
\textbf{读表提示}：1. Na 的 \(I_2\) 是 \(I_1\) 的 9.2 倍---第二个电子来自内层(2p)，这是判断 IA 族的标志。2. Al 的 \(I_3\to I_4\) 剧增 4.2 倍，说明失去 3 个电子后达到稳定构型\箭IIIA 族。3. Fe 的逐级 IE 无明显剧增(d 电子逐步失去)，这是过渡元素的典型特征。
\end{quote}

\subsubsection{2.3 电子亲和能}\label{ux7535ux5b50ux4eb2ux548cux80fd}

\textbf{定义}：基态气态原子获得一个电子成为一价气态负离子所放出的能量，称为第一电子亲和能 \(E_1\)。正值表示放热(通常情况)，负值表示吸热。

\textbf{第一电子亲和能数据}(kJ/mol)：

{\def\LTcaptype{none} % do not increment counter
\begin{longtable}[]{@{}cccccc@{}}
\toprule\noalign{}
元素 & \(E_1\) & 元素 & \(E_1\) & 元素 & \(E_1\) \\
\midrule\noalign{}
\endhead
\bottomrule\noalign{}
\endlastfoot
C & 122 & O & 141 & F & 328 \\
N & \约等于0 & S & 200 & Cl & \textbf{349} \\
P & 72 & Se & 195 & Br & 325 \\
\end{longtable}
}

\begin{quote}
\textbf{补充说明}：N 的 \(E_1 \approx 0\)---\(2\mathrm{p}^3\) 半满构型稳定，加电子需进入已有电子的 p 轨道，排斥大。类似地，稀有气体和 Be 的 \(E_1\) 也为负值或接近零。
\end{quote}

\textbf{第二电子亲和能恒为吸热}：负离子再获电子时，外来电子受负离子排斥，需吸收能量。例如 O\textsuperscript{-} + e\textsuperscript{-} \箭 O\textsuperscript{2-}，\(\Delta H = +780\) kJ/mol。因此 O\textsuperscript{2-}、S\textsuperscript{2-} 等二价负离子只能存在于晶体中(靠晶格能稳定)。

\textbf{周期趋势}

\textbf{同周期}(从左到右)：\(E_1\) \textbf{总体递增}(非金属性增强)。

\textbf{同主族}(从上到下)：\(E_1\) \textbf{总体递减}(半径增大，核对外来电子的引力减弱)。

\begin{figure}[H]\centering
\includegraphics[width=0.55\textwidth,keepaspectratio]{electron_affinity_trend.png}
\caption{元素第一电子亲和能随原子序数的变化。}
\end{figure}

\begin{quote}
\textbf{读图提示}：卤素(F、Cl、Br、I)的 EA 最大(放热最多)，稀有气体和IIA 族(ns\textsuperscript{2} 全满)的 EA 为负值或接近零(加电子需进入高能轨道，吸热)。注意 Cl 的峰值高于 F---这就是 §2.3 反常点的视觉证据。
\end{quote}

\textbf{关键反常点}

{\def\LTcaptype{none} % do not increment counter
\begin{longtable}[]{@{}
  >{\raggedright\arraybackslash}p{(\linewidth - 4\tabcolsep) * \real{0.3333}}
  >{\raggedright\arraybackslash}p{(\linewidth - 4\tabcolsep) * \real{0.3333}}
  >{\raggedright\arraybackslash}p{(\linewidth - 4\tabcolsep) * \real{0.3333}}@{}}
\toprule\noalign{}
\begin{minipage}[b]{\linewidth}\raggedright
反常
\end{minipage} & \begin{minipage}[b]{\linewidth}\raggedright
数据对比 (kJ/mol)
\end{minipage} & \begin{minipage}[b]{\linewidth}\raggedright
原因
\end{minipage} \\
\midrule\noalign{}
\endhead
\bottomrule\noalign{}
\endlastfoot
\textbf{Cl \textgreater{} F} & Cl(349) \textgreater{} F(328) & F 原子半径过小，电子云密度过高，额外引入电子时电子间排斥增大；Cl 有空 3d 轨道可缓冲 \\
\textbf{S \textgreater{} O} & S(200) \textgreater{} O(141) & 同上(O 的半径小，电子排斥大) \\
\end{longtable}
}

\begin{quote}
\textbf{易错提醒}：``非金属性越强，电子亲和能越大''不一定成立。F 是非金属性最强的元素，但电子亲和能 Cl \textgreater{} F。\textbf{EA 大 \不等于 单质反应活性高}---F\textsubscript{2} 的反应活性高于 Cl\textsubscript{2}，是因为 F\textsubscript{2} 的键能低(F-F 键弱)，而非 EA 大。反应活性取决于键能、晶格能、水合能等综合因素。
\end{quote}

\begin{quote}
\textbf{关联知识点}：元素周期律\#电子亲和能反常
\end{quote}

\subsubsection{2.4 电负性}\label{ux7535ux8d1fux6027}

\begin{quote}
\textbf{理解要点}：电负性是连接原子结构和分子结构的桥梁---从它可以判断键型(离子键vs共价键)、键的极性、甚至分子极性。没有电负性，你说不清为什么 HF 是极性分子而 F\textsubscript{2} 不是。
\end{quote}

\textbf{定义}：原子在分子中吸引成键电子对能力的相对量度。电负性越大的元素，非金属性越强；电负性越小的元素，金属性越强。

\textbf{主要标度}

{\def\LTcaptype{none} % do not increment counter
\begin{longtable}[]{@{}
  >{\raggedright\arraybackslash}p{(\linewidth - 6\tabcolsep) * \real{0.2500}}
  >{\raggedright\arraybackslash}p{(\linewidth - 6\tabcolsep) * \real{0.2500}}
  >{\raggedright\arraybackslash}p{(\linewidth - 6\tabcolsep) * \real{0.2500}}
  >{\raggedright\arraybackslash}p{(\linewidth - 6\tabcolsep) * \real{0.2500}}@{}}
\toprule\noalign{}
\begin{minipage}[b]{\linewidth}\raggedright
标度
\end{minipage} & \begin{minipage}[b]{\linewidth}\raggedright
提出者
\end{minipage} & \begin{minipage}[b]{\linewidth}\raggedright
方法
\end{minipage} & \begin{minipage}[b]{\linewidth}\raggedright
特征
\end{minipage} \\
\midrule\noalign{}
\endhead
\bottomrule\noalign{}
\endlastfoot
\textbf{Pauling 标度} \(\chi_p\) & Pauling (1932) & 由键能数据推算，规定 \(\chi_{\mathrm{F}}=3.98\) & 相对值，最常用 \\
\textbf{Mulliken 标度} \(\chi_M\) & Mulliken (1934) & \(\chi_M = \dfrac{1}{2}(I + A)\)(\(I\) 电离能，\(A\) 电子亲和能，单位 eV) & 绝对值，有明确物理意义 \\
\textbf{Allred-Rochow 标度} & Allred-Rochow (1957) & \(\chi_{AR} = 3590\dfrac{Z^*}{r^2} + 0.744\)(\(r\) 单位 pm) & 基于有效核电荷，与 Pauling 标度高度相关 \\
\end{longtable}
}

\begin{quote}
\textbf{补充说明}：\(\chi_p \approx 1.35\,\chi_M - 1.37\)(近似关系，两种标度的零点不同)
\end{quote}

\textbf{1. 算一算：用 Allred-Rochow 公式估算电负性}

已知 Na 的原子半径 186 pm，\(Z^* \approx 2.2\)：
\[\chi_{AR}(\mathrm{Na}) = 3590 \times \frac{2.2}{186^2} + 0.744 = 3590 \times \frac{2.2}{34596} + 0.744 \approx 0.23 + 0.744 \approx 0.97\]

Pauling 标度中 Na 的电负性为 0.93 --- 两种标度结果非常接近！

已知 Cl 的原子半径 99 pm，\(Z^* \approx 6.1\)，试试估算 Cl 的 \(\chi_{AR}\)：
\[\chi_{AR}(\mathrm{Cl}) = 3590 \times \frac{6.1}{99^2} + 0.744 = 3590 \times \frac{6.1}{9801} + 0.744 \approx 2.23 + 0.744 \approx 2.98\]

Pauling 标度中 Cl 的电负性为 3.16 --- 同样接近。这说明 \textbf{Allred-Rochow 标度的本质就是把''有效核电荷 ÷ 半径\textsuperscript{2}``映射到电负性}。

\begin{quote}
\textbf{点拨}：三种标度各有用途：Pauling 最常用(记忆数据用这个)，Mulliken 物理意义最清晰(就是电离能和电子亲和能的平均值)，Allred-Rochow 最''结构化学''(把电负性还原为 \(Z^*\) 和半径的关系)。其中 Allred-Rochow 公式最适合用来建立\textbf{电负性的物理直觉}---电负性大的本质是''核对外层电子的吸引力强''。
\end{quote}

\textbf{周期趋势}

\textbf{同周期}(从左到右)：电负性\textbf{递增}(稀有气体除外)。
\textbf{同主族}(从上到下)：电负性\textbf{递减}。

\textbf{规律}：F(3.98)电负性最大，Cs(0.79)最小。金属电负性一般 \textless{} 2.0，非金属一般 \textgreater{} 2.0。

\textbf{电负性的应用}

\begin{enumerate}
\def\labelenumi{\arabic{enumi}.}
\tightlist
\item
  \textbf{判断键型}：电负性差值 \(\Delta\chi < 1.7\) \箭 共价键为主；\(\Delta\chi > 1.7\) \箭 离子键为主(粗略估计)。\textbf{注意}：HF 的 \(\Delta\chi = 1.78 > 1.7\) 仍是强极性共价键---1.7 是经验阈值，不是严格分界
\item
  \textbf{判断氧化数}：化合物中电负性大的元素显负氧化态
\item
  \textbf{判断分子极性}：\(\Delta\chi\) 越大，键的极性越强
\end{enumerate}

\textbf{电负性受多种因素影响}：

{\def\LTcaptype{none} % do not increment counter
\begin{longtable}[]{@{}
  >{\raggedright\arraybackslash}p{(\linewidth - 4\tabcolsep) * \real{0.3333}}
  >{\raggedright\arraybackslash}p{(\linewidth - 4\tabcolsep) * \real{0.3333}}
  >{\raggedright\arraybackslash}p{(\linewidth - 4\tabcolsep) * \real{0.3333}}@{}}
\toprule\noalign{}
\begin{minipage}[b]{\linewidth}\raggedright
因素
\end{minipage} & \begin{minipage}[b]{\linewidth}\raggedright
规律
\end{minipage} & \begin{minipage}[b]{\linewidth}\raggedright
示例
\end{minipage} \\
\midrule\noalign{}
\endhead
\bottomrule\noalign{}
\endlastfoot
\textbf{氧化态} & 氧化态越高，电负性越大 & Fe\textsuperscript{3+}(1.96) \textgreater{} Fe\textsuperscript{2+}(1.83) \\
\textbf{杂化方式} & s 成分越高，电负性越大 & sp C(3.29) \textgreater{} sp\textsuperscript{2} C(2.75) \textgreater{} sp\textsuperscript{3} C(2.55) \\
\textbf{基团效应} & 吸电子基增大基团电负性 & −CF\textsubscript{3} \textgreater{} −CCl\textsubscript{3} \textgreater{} −CH\textsubscript{3} \\
\end{longtable}
}

\subsubsection{2.5 四大参数趋势总表}\label{ux56dbux5927ux53c2ux6570ux8d8bux52bfux603bux8868}

\begin{quote}
\textbf{关联知识点}：各项趋势的详细讨论见 §2.1-2.4，基础原理(Z*、屏蔽效应)见 原子结构-超级充实版(自学完整)\#3.6 周期律四趋势。
\end{quote}

{\def\LTcaptype{none} % do not increment counter
\begin{longtable}[]{@{}
  >{\raggedright\arraybackslash}p{(\linewidth - 8\tabcolsep) * \real{0.1739}}
  >{\centering\arraybackslash}p{(\linewidth - 8\tabcolsep) * \real{0.2174}}
  >{\centering\arraybackslash}p{(\linewidth - 8\tabcolsep) * \real{0.2174}}
  >{\centering\arraybackslash}p{(\linewidth - 8\tabcolsep) * \real{0.2174}}
  >{\raggedright\arraybackslash}p{(\linewidth - 8\tabcolsep) * \real{0.1739}}@{}}
\toprule\noalign{}
\begin{minipage}[b]{\linewidth}\raggedright
参数
\end{minipage} & \begin{minipage}[b]{\linewidth}\centering
符号
\end{minipage} & \begin{minipage}[b]{\linewidth}\centering
同周期(左\箭右)
\end{minipage} & \begin{minipage}[b]{\linewidth}\centering
同主族(上\箭下)
\end{minipage} & \begin{minipage}[b]{\linewidth}\raggedright
主要反常
\end{minipage} \\
\midrule\noalign{}
\endhead
\bottomrule\noalign{}
\endlastfoot
原子半径 & \(r\) & \下箭 减小 & \上箭 增大 & 稀有气体突增(范德华半径) \\
第一电离能 & \(I_1\) & \上箭 增大(整体) & \下箭 减小 & IIA \textgreater{} IIIA(\(s^2\) 全满)；VA \textgreater{} VIA(\(p^3\) 半满) \\
电子亲和能 & \(E_1\) & \上箭 增大(整体) & \下箭 减小 & F \textless{} Cl；O \textless{} S(半径过小\箭电子排斥) \\
电负性 & \(\chi\) & \上箭 增大 & \下箭 减小 & --- \\
\end{longtable}
}

\begin{center}\rule{0.5\linewidth}{0.5pt}\end{center}

\subsection{三、周期反常现象}\label{ux4e09ux5468ux671fux53cdux5e38ux73b0ux8c61}

\begin{quote}
\textbf{理解要点}：如果周期律只是''从左到右增大、从上到下增大''几个箭头，那它就不是化学了。真正的魅力在于那些''不听话''的地方---反常才是理解本质的窗口。每一个反常背后，都有电子排布的深层原因：惰性电子对效应告诉你相对论效应如何改写元素性质，对角线规则告诉你电荷密度比族更重要，第四周期非金属的最高价不稳定性告诉你过渡金属的存在如何影响隔壁p区元素。
\end{quote}

\subsubsection{3.1 对角规则}\label{ux5bf9ux89d2ux89c4ux5219}

第二周期元素与同族其他元素性质差异显著，却与第三周期\textbf{斜对角}元素性质相似，此现象称为\textbf{对角规则}。

\textbf{本质原因}：电荷密度相近(离子电荷/离子半径的比值相似)\箭 极化能力相似 \箭 化学性质相似。

\begin{quote}
\textbf{关联知识点}：元素周期律\#对角线规则 vs 常规族规律
\end{quote}

{\def\LTcaptype{none} % do not increment counter
\begin{longtable}[]{@{}
  >{\raggedright\arraybackslash}p{(\linewidth - 2\tabcolsep) * \real{0.5000}}
  >{\raggedright\arraybackslash}p{(\linewidth - 2\tabcolsep) * \real{0.5000}}@{}}
\toprule\noalign{}
\begin{minipage}[b]{\linewidth}\raggedright
对角对
\end{minipage} & \begin{minipage}[b]{\linewidth}\raggedright
相似性表现
\end{minipage} \\
\midrule\noalign{}
\endhead
\bottomrule\noalign{}
\endlastfoot
\textbf{Li - Mg} & 1. 过量 O\textsubscript{2} 中燃烧生成\textbf{氧化物}(非过氧化物)；2. 直接与 N\textsubscript{2} 反应生成氮化物(Li\textsubscript{3}N、Mg\textsubscript{3}N\textsubscript{2})；3. 氟化物/碳酸盐/磷酸盐难溶于水；4. 碳酸盐受热易分解；5. 易生成有机金属化合物 \\
\textbf{Be - Al} & 1. 金属/氧化物/氢氧化物均为\textbf{两性}；2. BeCl\textsubscript{2}、AlCl\textsubscript{3} 为缺电子共价化合物，易升华；3. Be、Al 遇冷浓硝酸钝化；4. 盐易水解 \\
\textbf{B - Si} & 1. 均不形成简单阳离子；2. 易形成可燃性氢化物(硼烷、硅烷)；3. 卤化物易水解 \\
\end{longtable}
}

\begin{quote}
\textbf{名师点拨}：对角规则的本质是\textbf{电荷密度相似}而非族相似。Li\textsuperscript{+}(Z=3, r=76pm) 的电荷密度与 Mg\textsuperscript{2+}(Z=12, r=72pm) 相近，因此极化能力相近，在形成共价性较强的化合物时行为相似。对比 Li\textsuperscript{+}(电荷密度\约等于0.52) 与同族的 Na\textsuperscript{+}(电荷密度\约等于0.18)，反而差异更大。
\end{quote}

\begin{quote}
\textbf{练一练}：电荷密度对比验证对角规则
\end{quote}

已知离子半径 Li\textsuperscript{+} 76 pm、Mg\textsuperscript{2+} 72 pm、Na\textsuperscript{+} 102 pm：
- Li\textsuperscript{+} 电荷密度 = 1 / (4\派×76\textsuperscript{2}/3) \约等于 0.52
- Mg\textsuperscript{2+} 电荷密度 = 2 / (4\派×72\textsuperscript{2}/3) \约等于 1.29
- Na\textsuperscript{+} 电荷密度 = 1 / (4\派×102\textsuperscript{2}/3) \约等于 0.18

Li\textsuperscript{+} 和 Mg\textsuperscript{2+} 的电荷密度确实比 Li\textsuperscript{+} 和 Na\textsuperscript{+} 更接近(虽不是严格相等，但在极化能力上已足够相似以产生可比拟的化学行为)。

\begin{quote}
\textbf{易错提醒}：``对角规则 = 性质完全相同''是错误的。对角规则只是\textbf{电荷密度相近导致的趋势相似}，并非等同。Li 和 Mg 在部分性质上相似，但在许多方面仍有本质差异(如 Li 的氢氧化物为弱碱，Mg(OH)\textsubscript{2} 为中强碱；Li 的氟化物溶解度远小于 MgF\textsubscript{2})。
\end{quote}

\begin{quote}
\textbf{竞赛常见考法}：以 Li 的异常性质为切入点，判断 Li 在周期表中的''反常''行为属于对角线规则，而非典型碱金属性质。
\end{quote}

\textbf{Li 的''反碱金属''行为}---Li 与 Mg 的对角相似导致它偏离典型碱金属性质：

{\def\LTcaptype{none} % do not increment counter
\begin{longtable}[]{@{}
  >{\raggedright\arraybackslash}p{(\linewidth - 6\tabcolsep) * \real{0.2500}}
  >{\raggedright\arraybackslash}p{(\linewidth - 6\tabcolsep) * \real{0.2500}}
  >{\raggedright\arraybackslash}p{(\linewidth - 6\tabcolsep) * \real{0.2500}}
  >{\raggedright\arraybackslash}p{(\linewidth - 6\tabcolsep) * \real{0.2500}}@{}}
\toprule\noalign{}
\begin{minipage}[b]{\linewidth}\raggedright
性质
\end{minipage} & \begin{minipage}[b]{\linewidth}\raggedright
Li(对角类 Mg)
\end{minipage} & \begin{minipage}[b]{\linewidth}\raggedright
Na/K(典型碱金属)
\end{minipage} & \begin{minipage}[b]{\linewidth}\raggedright
Mg
\end{minipage} \\
\midrule\noalign{}
\endhead
\bottomrule\noalign{}
\endlastfoot
氟化物溶解度 & LiF \textbf{难溶} & NaF/KF 可溶 & MgF\textsubscript{2} 难溶 \\
碳酸盐热稳定性 & Li\textsubscript{2}CO\textsubscript{3} \textbf{不稳定}，加热分解 & Na\textsubscript{2}CO\textsubscript{3} 极稳定 & MgCO\textsubscript{3} 不稳定 \\
硝酸盐分解 & LiNO\textsubscript{3} \箭 Li\textsubscript{2}O + NO\textsubscript{2} + O\textsubscript{2} & NaNO\textsubscript{3} \箭 NaNO\textsubscript{2} + O\textsubscript{2} & Mg(NO\textsubscript{3})\textsubscript{2} \箭 MgO + NO\textsubscript{2} + O\textsubscript{2} \\
\end{longtable}
}

\textbf{适用范围与限制}：对角规则\textbf{仅适用于前三周期主族}。第四周期及以下因 d 区插入和镧系收缩，对角相似性显著弱化。对角相似 \不等于 族相似---同族相似(价电子构型相同)强度远大于对角相似(电荷密度相近)。

\subsubsection{3.2 第四周期非金属元素最高价态不稳定性}\label{ux7b2cux56dbux5468ux671fux975eux91d1ux5c5eux5143ux7d20ux6700ux9ad8ux4ef7ux6001ux4e0dux7a33ux5b9aux6027}

第四周期非金属元素 As、Se、Br 的最高价态不稳定。Br 是最典型例证---高溴酸(HBrO\textsubscript{4})直到 1968 年才被制备出来(\$ \mathrm{NaBrO_3 + XeF_2 + H_2O \to NaBrO_4 + 2HF + Xe\uparrow} \$)。

\textbf{原因}：第一过渡系的存在使这些元素的有效核电荷增大，达最高价态所需激发能量超过总键能增加。

\begin{quote}
\textbf{名师点拨}：这一现象的本质是''过渡金属的副作用''：3d 电子的屏蔽不充分，使得第四周期 p 区元素的 Z* 比第三周期同族大得多。学生常见的思维盲区是孤立看每个元素，忽略了相邻过渡系对 p 区元素的''牵连影响''。
\end{quote}

\begin{quote}
\textbf{关联知识点}：元素周期律\#适用条件与限制(过渡金属的叠加效应)
\end{quote}

\subsubsection{3.3 惰性电子对效应与相对论效应}\label{ux60f0ux6027ux7535ux5b50ux5bf9ux6548ux5e94ux4e0eux76f8ux5bf9ux8bbaux6548ux5e94}

\textbf{定义}：第六周期 p 区元素(Tl、Pb、Bi)不易显示最高族价，而倾向于以比最高族价小 2 的氧化态存在。如 Tl(I) \textgreater{} Tl(III)、Pb(II) \textgreater{} Pb(IV)、Bi(III) \textgreater{} Bi(V)。

\textbf{本质}：\(6\mathrm{s}^2\) 电子对因\textbf{相对论效应}而收缩，钻穿效应强，平均能量低，不易参与成键。

\textbf{相对论效应机制}：重原子(如 Au、Hg、Pb)的内层电子运动速度可达光速的 58\%，相对论性质量增加：

\[m_{\text{rel}} = \frac{m_0}{\sqrt{1-(v/c)^2}}\]

质量增加 \箭 玻尔半径收缩(\(a_0 \propto 1/m\))\箭 **s 和 p 轨道向核收缩、能量降低**(直接效应)。s 电子收缩后对外层 d/f 电子的屏蔽增强 \箭 d/f 轨道能量相对升高、半径膨胀(间接效应)。

\textbf{结果}：

{\def\LTcaptype{none} % do not increment counter
\begin{longtable}[]{@{}
  >{\raggedright\arraybackslash}p{(\linewidth - 6\tabcolsep) * \real{0.2500}}
  >{\raggedright\arraybackslash}p{(\linewidth - 6\tabcolsep) * \real{0.2500}}
  >{\raggedright\arraybackslash}p{(\linewidth - 6\tabcolsep) * \real{0.2500}}
  >{\raggedright\arraybackslash}p{(\linewidth - 6\tabcolsep) * \real{0.2500}}@{}}
\toprule\noalign{}
\begin{minipage}[b]{\linewidth}\raggedright
元素
\end{minipage} & \begin{minipage}[b]{\linewidth}\raggedright
正常预期
\end{minipage} & \begin{minipage}[b]{\linewidth}\raggedright
实际性质
\end{minipage} & \begin{minipage}[b]{\linewidth}\raggedright
相对论效应解释
\end{minipage} \\
\midrule\noalign{}
\endhead
\bottomrule\noalign{}
\endlastfoot
\textbf{Au} & 银白色(如 Cu、Ag) & \textbf{金黄色} & 6s\箭5d 跃迁能量落入可见光区(蓝紫光被吸收) \\
\textbf{Hg} & 常温固态金属 & \textbf{液态}(熔点 -39°C) & 6s\textsuperscript{2} 收缩导致金属键减弱(6s 电子不易参与金属键) \\
\textbf{Tl/Pb} & 最高价稳定 & \textbf{低价更稳定} & 6s\textsuperscript{2} 电子能量低，不易参与成键 \箭 惰性电子对效应 \\
\end{longtable}
}

这一效应在 \(n\mathrm{s}^2\) 中随 \(n\) 增大而增强(\(6\mathrm{s}^2\) 最显著)，是竞赛中解释 Pb(II) 稳定性、Au 颜色和 Hg 液态的核心依据。

\begin{quote}
\textbf{点拨}：这里有一个经典的认知冲突：同族从上到下，金属性增强，按理说\textbf{低价态越来越稳定}(如 C(IV)\textgreater C(II)、Si(IV)\textgreater Si(II)、Ge(IV)\textgreater Ge(II))，但到了第六周期突然\textbf{反转}：Pb(II) \textgreater{} Pb(IV)！容易认为''同族往下都是高价更稳定''，但惰性电子对效应告诉：\textbf{相对论效应可以改写周期趋势}。
\end{quote}

对比同族价态稳定性：
- C: IV \textgreater\textgreater{} II
- Si: IV \textgreater\textgreater{} II
- Ge: IV \textgreater{} II
- Sn: IV \大于等于 II
- \textbf{Pb: II \textgreater{} IV} \左箭 临界反转点！

\begin{quote}
\textbf{练一练}：用惰性电子对效应解释异常现象
\end{quote}

\textbf{问题}：已知 Tl(I) 的化合物比 Tl(III) 更稳定，而 Al(III) 的化合物比 Al(I) 更稳定。请用惰性电子对效应解释。

\textbf{思路}：Al 在第三周期(\(n=3\))，\(3\mathrm{s}^2\) 电子对的钻穿效应和相对论收缩不明显，容易参与成键 \箭 Al(III) 稳定。Tl 在第六周期(\(n=6\))，\(6\mathrm{s}^2\) 电子对因相对论效应显著收缩 \箭 钻穿效应强 \箭 6s 电子能量低 \箭 ``不想''参与成键 \箭 Tl(I) 稳定。

\begin{quote}
\textbf{深化理解}：惰性电子对效应不是孤立现象，它是\textbf{相对论效应}在周期表中的重要表现之一(另两个是金的颜色和汞的液态)。这是为什么竞赛中常把 Pb(II) 稳定性、Au 的惰性、Hg 的低熔点放在一起考。
\end{quote}

\textbf{标准电极电势的定量证据}：

{\def\LTcaptype{none} % do not increment counter
\begin{longtable}[]{@{}lcl@{}}
\toprule\noalign{}
电对 & \(E^\circ\) (V) & 含义 \\
\midrule\noalign{}
\endhead
\bottomrule\noalign{}
\endlastfoot
Sn\textsuperscript{4+}/Sn\textsuperscript{2+} & +0.15 & Sn(IV) 氧化性弱 \\
Pb\textsuperscript{4+}/Pb\textsuperscript{2+} & \textbf{+1.46} & Pb(IV) 是强氧化剂 \\
\end{longtable}
}

从 Sn 到 Pb，高氧化态的稳定性急剧下降---这就是惰性电子对效应的电化学定量表现。

\textbf{竞赛常见应用}：
- \textbf{铅蓄电池正极}：PbO\textsubscript{2} + Pb + 2H\textsubscript{2}SO\textsubscript{4} \箭 2PbSO\textsubscript{4} + 2H\textsubscript{2}O(Pb\textsuperscript{4+} \箭 Pb\textsuperscript{2+} 还原)
- \textbf{NaBiO\textsubscript{3} 定性检验 Mn\textsuperscript{2+}}：NaBiO\textsubscript{3} 可将 Mn\textsuperscript{2+} 氧化为 MnO\textsubscript{4}\textsuperscript{-}(紫红色)，是 Mn\textsuperscript{2+} 的特征鉴定反应
- \textbf{Pb(IV) 的''复活''}：强配体如 {[}PbCl\textsubscript{6}{]}\textsuperscript{2-} 可稳定 Pb(IV)，说明惰性对效应是热力学倾向，可被强配位场克服

\subsubsection{3.4 氢的特殊性}\label{ux6c22ux7684ux7279ux6b8aux6027}

氢在周期表中位置特殊，兼具多种族的特征：

{\def\LTcaptype{none} % do not increment counter
\begin{longtable}[]{@{}
  >{\raggedright\arraybackslash}p{(\linewidth - 4\tabcolsep) * \real{0.3333}}
  >{\raggedright\arraybackslash}p{(\linewidth - 4\tabcolsep) * \real{0.3333}}
  >{\raggedright\arraybackslash}p{(\linewidth - 4\tabcolsep) * \real{0.3333}}@{}}
\toprule\noalign{}
\begin{minipage}[b]{\linewidth}\raggedright
类比对象
\end{minipage} & \begin{minipage}[b]{\linewidth}\raggedright
依据
\end{minipage} & \begin{minipage}[b]{\linewidth}\raggedright
差异
\end{minipage} \\
\midrule\noalign{}
\endhead
\bottomrule\noalign{}
\endlastfoot
\textbf{IA 族} & \(1\mathrm{s}^1\) 最外层 1 个电子 \箭 可失电子成 H\textsuperscript{+} & H\textsuperscript{+} 半径极小(\(1.5\times10^{-15}\) m)，性质与碱金属离子差异极大 \\
\textbf{VIIA 族} & 获得 1 个电子成 H\textsuperscript{-} \箭 稀有气体构型 & H\textsuperscript{-} 半径大(154 pm)，还原性强；电负性(2.2) 远小于卤素 \\
\textbf{IVA 族} & 价电子半满 \箭 既可作氧化剂又可作还原剂 & --- \\
\end{longtable}
}

\begin{quote}
\textbf{点拨}：氢的''三不像''特性本身就是最好的素材：周期表的位置不是绝对的，而是由电子构型带来的\textbf{倾向性}决定的。竞赛中常问''为什么 H 放在 IA 族上方''，标准答案不是''因为它像碱金属''，而是因为它有 1 个 s 电子，电子构型与IA一致''**。
\end{quote}

\subsubsection{3.5 常见误区}\label{ux5e38ux89c1ux8befux533a}

\begin{quote}
\textbf{误区1}：``同周期从左到右电离能单调递增''---实际有IIA/IIIA、VA/VIA 两处反常，根源是电子排布的稳定性(\(s^2\) 全满、\(p^3\) 半满额外稳定)。
\end{quote}

\begin{quote}
\textbf{误区2}：``电离能趋势与原子半径趋势完全相反''---大致相反但不是严格一一对应。电离能还受电子排布稳定性的影响(如 N \textgreater{} O 是半满稳定因素 \textgreater{} 半径因素)。
\end{quote}

\begin{quote}
\textbf{误区3}：``镧系收缩仅影响镧系元素''---它的主要影响反而是使第二、第三过渡系同族元素半径相近、性质相似(Zr/Hf、Nb/Ta、Mo/W)。这是竞赛高频考点。
\end{quote}

\begin{quote}
\textbf{误区4}：``对角规则 = 性质完全相同''---只是电荷密度相近导致的趋势相似，并非等同。Li 和 Mg 在部分性质上相似，但在许多方面仍有差异(详见 §3.1)。
\end{quote}

\begin{quote}
\textbf{误区5}：``电子亲和能越大的元素非金属性越强''---大体正确但 Cl \textgreater{} F 是反例。F 的非金属性最强，但电子亲和能 Cl \textgreater{} F(源于 F 半径过小\箭电子排斥)。
\end{quote}

\begin{quote}
\textbf{误区6}：``电负性差值 \textgreater{} 1.7 一定是离子键''---这是粗略估计而非严格判据。如 HF 的 \(\Delta\chi=1.78>1.7\) 却是强极性共价键。键型判断除了看\(\Delta\chi\)，还要考虑元素种类和分子结构。
\end{quote}

\begin{quote}
\textbf{口诀}：
- ``IIA IIIA 绊一跤，VA VIA 也摔倒''---电离能两处反常
- ``镧系一收缩，Zr/Hf 分不清''---镧系收缩的化学后果
\end{quote}

\subsubsection{3.6 次级周期性---过渡金属的''副作用''}\label{ux6b21ux7ea7ux5468ux671fux6027ux8fc7ux6e21ux91d1ux5c5eux7684ux526fux4f5cux7528}

同族元素从上到下，性质并非单调变化。第四周期和第六周期分别因 \textbf{d 电子}和 \textbf{f 电子}的首次出现而产生''反常''---这就是\textbf{次级周期性}。

\textbf{因果链}：

\begin{verbatim}
第四周期：首次出现 d 电子
第六周期：首次出现 f 电子 + 相对论效应
        ↓
d/f 电子对外层电子屏蔽差 → 有效核电荷 Z* 比预期大
        ↓
外层电子更难电离 → 电离能偏高 → 外层电子"惰性"
        ↓
高价化合物倾向于夺回电子 → 氧化性增强
\end{verbatim}

\begin{quote}
\textbf{竞赛例证}：
- \textbf{第四周期效应}：浓热 \(\mathrm{H_2SeO_4}\) 可氧化 Au，而同族的 \(\mathrm{H_2SO_4}\) 不行(3d 电子屏蔽差 \箭 Se 的 Z* 偏大 \箭 Se(VI) 氧化性异常强)
- \textbf{第六周期效应}：\(\mathrm{PbO_2}\)、\(\mathrm{NaBiO_3}\) 可将 \(\mathrm{Mn^{2+}}\) 氧化为 \(\mathrm{MnO_4^-}\)(4f 电子屏蔽更差 + 相对论效应 \箭 Pb(IV)、Bi(V) 强氧化性)
- \textbf{惰性电子对}：Tl(I) 比 Tl(III) 稳定、Hg 为液态金属---都是第六周期 f 电子 + 相对论效应的叠加结果
\end{quote}

\begin{quote}
\textbf{关联知识点}：元素周期律\#次级周期性 · 相对论效应 · 惰性对效应
\end{quote}

\begin{center}\rule{0.5\linewidth}{0.5pt}\end{center}

\subsection{四、典型例题}\label{ux56dbux5178ux578bux4f8bux9898}

\textbf{题目}：比较下列各组原子/离子的半径大小，并说明判断依据：
(1) Na 与 Mg；(2) Na\textsuperscript{+} 与 F\textsuperscript{-}；(3) Fe\textsuperscript{2+} 与 Fe\textsuperscript{3+}；(4) Zr 与 Hf

\textbf{题目}：解释以下电离能数据反常的原因：
(1) 第一电离能 Be(900) \textgreater{} B(801)
(2) 第一电离能 N(1402) \textgreater{} O(1314)
(3) 第二电离能 O\textsuperscript{+}(2666) \textgreater{} F\textsuperscript{+}(2674)

\textbf{题目}：元素 A 的 \(I_1\sim I_5\) (kJ/mol) 分别为：578, 1817, 2745, 11575, 14830。推断 A 的族属，并写出其基态电子排布式。

\textbf{题目}：已知某元素基态原子的电子排布式为 \([\mathrm{Xe}] 4\mathrm{f}^{14} 5\mathrm{d}^{10} 6\mathrm{s}^1\)。判断：(1) 原子序数；(2) 周期、族、区；(3) 未成对电子数。

\textbf{题目}：解释以下现象：
(1) Zr 和 Hf 在自然界中总是共生，分离困难。
(2) 第三过渡系元素的电离能普遍高于第二过渡系同族元素。

\textbf{题目}：碳族元素(IVA)的常见氧化态变化规律为：C(IV) \textgreater{} C(II)，Si(IV) \textgreater{} Si(II)，Ge(IV) \textgreater{} Ge(II)，Sn(IV) \大于等于 Sn(II)，Pb(II) \textgreater{} Pb(IV)。请回答：
(1) 这一变化规律的名称是什么？
(2) 为什么到了 Pb 出现''反转''---Pb(II) 反而比 Pb(IV) 更稳定？
(3) 已知 Tl(IIIA 族)也有类似现象：Tl(I) 比 Tl(III) 更稳定。Bi(VA 族)则表现为 Bi(III) \textgreater{} Bi(V)。这些现象有何共同本质？

\begin{center}\rule{0.5\linewidth}{0.5pt}\end{center}

\textbf{例题答案}

\textbf{例1}：(1) Na \textgreater{} Mg(同周期左\箭右 \(Z^*\)\上箭 \(r\)\下箭)；(2) Na\textsuperscript{+} \textless{} F\textsuperscript{-}(同电子层，核电荷大者半径小)；(3) Fe\textsuperscript{2+} \textgreater{} Fe\textsuperscript{3+}(正电荷越高半径越小)；(4) Zr \约等于 Hf(镧系收缩使半径几乎相等)。

\textbf{例2}：(1) Be 的 \(2\mathrm{s}^2\) 全满稳定，B 失 \(2\mathrm{p}\) 后得 \(2\mathrm{s}^2\) 稳定构型；(2) N 的 \(2\mathrm{p}^3\) 半满稳定，O 失一个电子后达半满；(3) O\textsuperscript{+} 为 \(2\mathrm{p}^3\) 半满，F\textsuperscript{+} 为 \(2\mathrm{p}^4\) 非半满。

\textbf{例3}：\(I_3 \to I_4\) 剧增至 11575(4.2×)\箭 3 个价电子 \箭 IIIA 族 \箭 Al(\([\mathrm{Ne}]3\mathrm{s}^2 3\mathrm{p}^1\))。

\textbf{例4}：(1) Z=79 \箭 Au；(2) 第六周期 IB 族 ds 区；(3) 1 个未成对电子(\(6\mathrm{s}^1\))。

\textbf{例5}：(1) 镧系收缩使 Hf(159 pm) \约等于 Zr(160 pm)，化学性质极相似 \箭 分离困难；(2) 半径相近但核电荷大增 \箭 \(Z^*\) 更大 \箭 电离能更高。

\textbf{例6}：(1) 惰性电子对效应；(2) Pb 的 \(6\mathrm{s}^2\) 因相对论效应显著收缩，不易参与成键 \箭 Pb(II) 稳定；(3) 共同本质：第六周期 \(n\mathrm{s}^2\) 电子对的相对论稳定化，随 \(n\) 增大而增强。

\subsection{五、本节总结 · 速查}\label{ux4e94ux672cux8282ux603bux7ed3-ux901fux67e5}

{\def\LTcaptype{none} % do not increment counter
\begin{longtable}[]{@{}
  >{\raggedright\arraybackslash}p{(\linewidth - 4\tabcolsep) * \real{0.3333}}
  >{\raggedright\arraybackslash}p{(\linewidth - 4\tabcolsep) * \real{0.3333}}
  >{\raggedright\arraybackslash}p{(\linewidth - 4\tabcolsep) * \real{0.3333}}@{}}
\toprule\noalign{}
\begin{minipage}[b]{\linewidth}\raggedright
核心概念
\end{minipage} & \begin{minipage}[b]{\linewidth}\raggedright
关键规律
\end{minipage} & \begin{minipage}[b]{\linewidth}\raggedright
常见陷阱
\end{minipage} \\
\midrule\noalign{}
\endhead
\bottomrule\noalign{}
\endlastfoot
\textbf{周期表结构} & 周期 = 能级组，族 = 价电子数，区 = 最后填入能级 & 过渡元素价电子 = \((n-1)\mathrm{d} + n\mathrm{s}\)，非仅最外层 \\
\textbf{原子半径} & 同周期 \下箭(\(Z^*\)\上箭)，同族 \上箭(\(n\)\上箭) & 稀有气体用范德华半径，不可与共价半径直接比 \\
\textbf{电离能} & 同周期\上箭(两处反常)，同族\下箭 & IIA \textgreater{} IIIA(\(s^2\) 全满)，VA \textgreater{} VIA(\(p^3\) 半满) \\
\textbf{电子亲和能} & 同周期\上箭，同族\下箭 & Cl \textgreater{} F 是反例(F 过小\箭电子排斥) \\
\textbf{电负性} & F 最大 (3.98)，Cs 最小 (0.79) & 同一元素不同氧化态电负性不同；\(\Delta\chi>1.7\) 不一定是离子键 \\
\textbf{镧系收缩} & 镧系 15 种元素半径逐渐减小 & 主要影响是第二、第三过渡系同族相似 \\
\textbf{对角规则} & Li-Mg, Be-Al, B-Si & 性质相似 \不等于 等同，仅电荷密度相近 \\
\textbf{惰性电子对效应} & 第六周期 p 区低价稳定 & Tl(I)\textgreater Tl(III); Pb(II)\textgreater Pb(IV) \\
\textbf{交换能} & 平行自旋对越多 \箭 交换能越负 \箭 越稳定 & Cr(3d\textsuperscript{5}4s\textsuperscript{1}) 和 Cu(3d\textsuperscript{10}4s\textsuperscript{1}) 的异常排布 \\
\textbf{次级周期性} & 第四/六周期高价化合物氧化性增强 & d/f 电子屏蔽差 \箭 Z* 偏大 \箭 外层电子''惰性'' \\
\textbf{相对论效应} & 6s 轨道收缩 \箭 Au 颜色、Hg 液态、惰性电子对 & 不只是''边角知识''，是重元素反常的统一解释 \\
\textbf{元素推断} & 电离能突变点 \箭 价电子数 \箭 族属 & 注意过渡元素的价电子包括 d 电子 \\
\end{longtable}
}

\begin{center}\rule{0.5\linewidth}{0.5pt}\end{center}

\subsection{六、练习题}\label{ux516dux7ec3ux4e60ux9898}

\begin{enumerate}
\def\labelenumi{\arabic{enumi}.}
\item
  \textbf{判断下列元素位于第几周期、第几族、什么区}：

  \begin{enumerate}
  \def\labelenumii{(\arabic{enumii})}
  \tightlist
  \item
    Cl (Z=17)；(2) Fe (Z=26)；(3) Zn (Z=30)；(4) Br (Z=35)；(5) Cs (Z=55)
  \end{enumerate}
\item
  \textbf{价层电子构型与电离能}：下列不同原子的价层电子构型中，哪一个第一电离能最大？哪一个最小？

  \begin{enumerate}
  \def\labelenumii{(\arabic{enumii})}
  \tightlist
  \item
    \(2\mathrm{s}^2 2\mathrm{p}^1\)　(2) \(2\mathrm{s}^2 2\mathrm{p}^2\)　(3) \(2\mathrm{s}^2 2\mathrm{p}^3\)　(4) \(2\mathrm{s}^2 2\mathrm{p}^4\)
  \end{enumerate}
\item
  \textbf{碱金属性质递变}(第30届初赛)：对碱金属 Li、Na、K、Rb 和 Cs，随着原子序数增加，以下哪种性质的递变\textbf{不是单调的}？简述原因。

  \begin{enumerate}
  \def\labelenumii{(\alph{enumii})}
  \tightlist
  \item
    熔、沸点　(b) 原子半径　(c) 晶体密度　(d) 第一电离能
    \textgreater{} 来源：第30届中国化学奥林匹克(初赛)
  \end{enumerate}
\item
  \textbf{镧系收缩与离子半径}(第37届决赛)：依照元素物理和化学性质的特点，以下三种元素离子半径最大的是？
  A. Sc　B. Y　C. Lu
  \textgreater{} 来源：第37届中国化学奥林匹克(决赛)
\item
  \textbf{全面半径比较}(普化教材习题 11.19)：据理判定下列各对原子(或离子)中哪一个半径大，并简述理由。
  H 与 He；Ba 与 Sr；Sc 与 Ca；Cu 与 Ni；Zr 与 Hf；La 与 Gd；S\textsuperscript{2-} 与 S；Na 与 Al\textsuperscript{3+}；Fe\textsuperscript{2+} 与 Fe\textsuperscript{3+}；Pb\textsuperscript{2+} 与 Sn\textsuperscript{2+}
\item
  \textbf{等电子体半径差异}(普化教材习题 11.20)：试说明下列等电子体的离子半径值为什么在数值上有差别：
  F\textsuperscript{-}(133 pm)、O\textsuperscript{2-}(140 pm)、Na\textsuperscript{+}(102 pm)、Mg\textsuperscript{2+}(72 pm)、Al\textsuperscript{3+}(54 pm)
\item
  \textbf{反常解释}：

  \begin{enumerate}
  \def\labelenumii{(\arabic{enumii})}
  \tightlist
  \item
    解释为什么 Mg 的第一电离能大于 Al。
  \item
    解释为什么 P 的第一电离能大于 S。
  \item
    解释为什么 F 的电负性大于 Cl，但电子亲和能 Cl 大于 F。
  \end{enumerate}
\item
  \textbf{元素推断}：某元素的 \(I_1\sim I_5\) 数据如下 (kJ/mol)：
  \[I_1=786,\ I_2=1580,\ I_3=3230,\ I_4=4360,\ I_5=16000\]

  \begin{enumerate}
  \def\labelenumii{(\arabic{enumii})}
  \tightlist
  \item
    判断该元素有几个价电子，位于第几族。
  \item
    若该元素为第三周期元素，写出其元素符号和电子排布式。
  \item
    写出该元素最高价氧化物的水化物化学式，判断其酸碱性。
  \end{enumerate}
\item
  \textbf{镧系收缩应用}：已知 Zr 的原子半径为 160 pm，Hf 为 159 pm，而根据同族趋势规律，Hf(第六周期)的原子半径应大于 Zr(第五周期)。请解释为什么 Hf 的半径反而略小于 Zr，并说明这一现象对 Zr/Hf 分离产生什么影响。
\item
  \textbf{对角线规则应用}(无机化学教材习题)：

  \begin{enumerate}
  \def\labelenumii{(\arabic{enumii})}
  \tightlist
  \item
    列举 Li 与 Mg 性质相似的三个实例。
  \item
    列举 Be 与 Al 性质相似的三个实例。
  \item
    解释对角线规则的物理化学根源。
  \end{enumerate}
\item
  \textbf{综合推断}：A、B、C、D 为四种短周期元素。已知：

  \begin{itemize}
  \tightlist
  \item
    A 的简单阴离子与 D 的简单阳离子具有相同的电子层结构
  \item
    B 的最高价氧化物对应水化物与 D 的最高价氧化物对应水化物反应生成盐和水
  \item
    C 的单质既能与强酸反应又能与强碱反应
  \item
    A 的第一电离能在同周期元素中仅低于稀有气体
  \end{itemize}

  请推断 A、B、C、D 的元素符号，并说明推理过程。
\item
  \textbf{周期律综合分析}：X、Y、Z、W 是原子序数依次增大的短周期元素。已知：

  \begin{itemize}
  \tightlist
  \item
    X 的原子半径在同周期中最小
  \item
    Y 的简单离子与 W 的简单离子具有相同的电子层结构
  \item
    Z 的第一电离能在其同周期元素中最小
  \item
    W 的氧化物对应水化物为强酸
  \end{itemize}

  回答：(1) 写出 X、Y、Z、W 的元素符号；(2) 比较 X、Y、Z 的原子半径大小；(3) 比较 Y 和 W 的简单离子的半径大小；(4) 写出 X 和 Z 形成的化合物的化学式并判断键型。
\item
  \textbf{惰性电子对效应解释}：已知 Tl 在第六周期IIIA 族，其常见氧化态为 Tl(I) 和 Tl(III)，且 Tl(I) 比 Tl(III) 更稳定。请回答：

  \begin{enumerate}
  \def\labelenumii{(\arabic{enumii})}
  \tightlist
  \item
    这一现象叫什么？从电子结构角度解释原因。
  \item
    对比 Al(第三周期IIIA 族)的常见氧化态，说明为什么同为IIIA 族，趋势却不同。
  \item
    预测 In(第五周期IIIA 族)的 Tl(I)/Tl(III) 稳定性关系，说明理由。
  \end{enumerate}
\item
  \textbf{Bent 规则与电负性}(第39届初赛)：氟在元素周期表中具有最大的电负性。与 C-H、C-O、C-N 以及 C-Cl 等键相比，C-F 键具有最大的键解离能和最大的极性。

  \begin{enumerate}
  \def\labelenumii{(\arabic{enumii})}
  \tightlist
  \item
    比较甲烷和一氟甲烷中的 H-C-H 键角，以下说法准确的是(　)：
  \end{enumerate}

  \begin{enumerate}
  \def\labelenumii{(\alph{enumii})}
  \tightlist
  \item
    F 的电负性使 C-F 共价电子偏向 F，C-H 键之间排斥力减少，键角变大
  \item
    F 的电负性使 C-F 共价电子偏向 F，C-H 键之间排斥力增加，键角变小
  \item
    C-H 键之间排斥力没有变化，键角不变
  \end{enumerate}

  \begin{enumerate}
  \def\labelenumii{(\arabic{enumii})}
  \setcounter{enumii}{1}
  \tightlist
  \item
    与一氟甲烷相比，二氟甲烷中 H-C-H 键角如何变化？(a) 变大　(b) 变小　(c) 不变
    \textgreater{} 来源：第39届中国化学奥林匹克(初赛)
  \end{enumerate}
\end{enumerate}

\begin{center}\rule{0.5\linewidth}{0.5pt}\end{center}

\subsection{七、参考答案与提示}\label{ux4e03ux53c2ux8003ux7b54ux6848ux4e0eux63d0ux793a}

\textbf{6.1 基础巩固}

\begin{enumerate}
\def\labelenumi{\arabic{enumi}.}
\item
  \begin{enumerate}
  \def\labelenumii{(\arabic{enumii})}
  \tightlist
  \item
    Cl：第三周期 VIIA 族 p 区；(2) Fe：第四周期 VIII族 d 区；(3) Zn：第四周期 IIB 族 ds 区；(4) Br：第四周期 VIIA 族 p 区；(5) Cs：第六周期 IA 族 s 区。
  \end{enumerate}
\item
  I\textsubscript{1} 最大：(3) \(2\mathrm{s}^2 2\mathrm{p}^3\)(p\textsuperscript{3} 半满稳定，交换能最大)；I\textsubscript{1} 最小：(1) \(2\mathrm{s}^2 2\mathrm{p}^1\)(p\textsuperscript{1} 失去后得 s\textsuperscript{2} 全满稳定构型，且核电荷最小)。
\item
  \begin{enumerate}
  \def\labelenumii{(\alph{enumii})}
  \setcounter{enumii}{2}
  \tightlist
  \item
    晶体密度\textbf{不是单调的}---Na(0.97 g/cm\textsuperscript{3}) \textgreater{} K(0.89 g/cm\textsuperscript{3})。原因：碱金属为 BCC 结构，密度 = 原子质量 / 原子体积。从 Na 到 K，原子质量增加约 65\%，但原子体积增大约 170\%(半径增大导致)，体积增幅 \textgreater{} 质量增幅 \箭 密度反而下降。熔沸点、原子半径、第一电离能均为单调递变。
  \end{enumerate}
\item
  \textbf{B(Y)}。Sc\textsuperscript{3+}(74.5 pm) \textgreater{} Y\textsuperscript{3+}(90.0 pm) \textgreater{} Lu\textsuperscript{3+}(86.1 pm)。Y\textsuperscript{3+} 半径最大---Y 在第五周期，电子层数比 Sc 多一层；Lu 虽在第六周期，但镧系收缩使 Lu\textsuperscript{3+} 半径回缩到与 Y\textsuperscript{3+} 接近甚至更小。
\end{enumerate}

\textbf{6.2 提高训练}

\begin{enumerate}
\def\labelenumi{\arabic{enumi}.}
\setcounter{enumi}{4}
\item
  \textbf{提示}(10 对半径比较的核心思路)：

  \begin{itemize}
  \tightlist
  \item
    H \textless{} He(He 的 Z 更大，电子更紧缩)
  \item
    Ba \textgreater{} Sr(同族从上到下递增)
  \item
    Sc \textless{} Ca(同周期从左到右递减，且 Sc 的 Z* 更大)
  \item
    Cu \textgreater{} Ni(Cu 的 3d\textsuperscript{10} 全满，d 电子屏蔽差 \箭 有效半径偏大)
  \item
    Zr \约等于 Hf(镧系收缩)
  \item
    La \textgreater{} Gd(镧系收缩，La 在 Gd 前面，半径更大)
  \item
    S\textsuperscript{2-} \textgreater{} S(负离子 \textgreater{} 原子 \textgreater{} 正离子)
  \item
    Na \textgreater{} Al\textsuperscript{3+}(等电子体，核电荷大者半径小)
  \item
    Fe\textsuperscript{2+} \textgreater{} Fe\textsuperscript{3+}(正电荷越高半径越小)
  \item
    Pb\textsuperscript{2+} \textgreater{} Sn\textsuperscript{2+}(Pb\textsuperscript{2+} 有 4f 电子但镧系收缩影响有限，主要看电子层数)
  \end{itemize}
\item
  \textbf{提示}：五个离子电子层结构相同(均为 1s\textsuperscript{2}2s\textsuperscript{2}2p\textsuperscript{6})，但核电荷不同：O\textsuperscript{2-}(8) \textless{} F\textsuperscript{-}(9) \textless{} Na\textsuperscript{+}(11) \textless{} Mg\textsuperscript{2+}(12) \textless{} Al\textsuperscript{3+}(13)。核电荷越大 \箭 对相同数目的电子吸引力越强 \箭 半径越小。这是 Z* 对离子半径影响的最直观体现。
\item
  \textbf{提示}：(1) Mg 的 \(3\mathrm{s}^2\) 全满稳定；(2) P 的 \(3\mathrm{p}^3\) 半满稳定；(3) F 半径太小，增加电子时电子间排斥大 \箭 放热少；但电负性是衡量原子在分子中吸引成键电子的能力，F 的核吸引能力更强。
\item
  \textbf{提示}：(1) \(I_4\to I_5\) 剧增 \箭 4 个价电子 \箭 IVA 族。(2) 第三周期 IVA \箭 Si(\([\mathrm{Ne}] 3\mathrm{s}^2 3\mathrm{p}^2\))。(3) H\textsubscript{2}SiO\textsubscript{3}(或 H\textsubscript{4}SiO\textsubscript{4})，弱酸性。
\item
  \textbf{提示}：镧系收缩使第三过渡系元素(Hf)的半径反常减小，与第二过渡系同族(Zr)半径相近 \箭 化学性质相似 \箭 分离极为困难(通常用离子交换或溶剂萃取法分离)。
\item
  \textbf{提示}：(1) Li-Mg：1. 过量 O\textsubscript{2} 中燃烧生成氧化物(非过氧化物)；2. 直接与 N\textsubscript{2} 反应生成氮化物；3. 碳酸盐受热易分解。(2) Be-Al：1. 金属/氧化物/氢氧化物均为两性；2. 氯化物为缺电子共价化合物，易升华；3. 遇冷浓硝酸钝化。(3) 物理化学根源：对角元素的\textbf{离子势}(\(\phi = z/r\))相近 \箭 极化能力相近 \箭 化学性质相似。
\end{enumerate}

\textbf{6.3 挑战题}

\begin{enumerate}
\def\labelenumi{\arabic{enumi}.}
\setcounter{enumi}{10}
\item
  \textbf{提示}：A 的第一电离能仅次于稀有气体 \箭 A 为卤素(VIIA 族)。C 既能与强酸又能与强碱反应 \箭 C 为 Be 或 Al(短周期两性元素)。结合 A 为阴离子、D 为阳离子且电子层结构相同 \箭 A 为 F、C 为 Al、D 为 Na、B 为 Cl(NaOH + HClO\textsubscript{4} 反应)。答案：A=F，B=Cl，C=Al，D=Na。
\item
  \textbf{提示}：X 的原子半径同周期最小 \箭 卤素；Z 的第一电离能同周期最小 \箭 碱金属；Y 与 W 的离子电子层结构相同且 W 的氧化物对应水化物为强酸 \箭 W 为 Cl(或 S)。综合分析：X=F，Y=Na，Z=Mg，W=Cl(或 X=Cl，Y=Li，Z=Na，W=S)。注意不同情况的合理性。
\item
  \textbf{提示}：

  \begin{enumerate}
  \def\labelenumii{(\arabic{enumii})}
  \tightlist
  \item
    \textbf{惰性电子对效应}。Tl 的 \(6\mathrm{s}^2\) 电子对因相对论效应显著收缩 \箭 钻穿效应强 \箭 6s 电子能量低 \箭 不易参与成键 \箭 Tl(I) 保留 \(6\mathrm{s}^2\) 更稳定。
  \item
    Al 在第三周期(\(n=3\))，\(3\mathrm{s}^2\) 电子对的相对论效应和钻穿效应不明显 \箭 容易参与成键 \箭 Al(III) 稳定。同为 IIIA 族但因 n 值不同导致 s 电子对的惰性程度不同。
  \item
    In 在第五周期，\(5\mathrm{s}^2\) 有一定惰性电子对效应但远弱于 Tl 的 \(6\mathrm{s}^2\) \箭 In(III) 仍比 In(I) 稳定，但 In(I) 的稳定性已高于 Al(I) 和 Ga(I)。趋势：In(III) \textgreater{} In(I)，但差距缩小。
  \end{enumerate}
\item
  \textbf{提示}：

  \begin{enumerate}
  \def\labelenumii{(\arabic{enumii})}
  \tightlist
  \item
    选 \textbf{(a)}。根据 \textbf{Bent 规则}：电负性大的取代基(F)倾向于占据 p 成分更多的杂化轨道 \箭 C-F 键使用更多 p 成分 \箭 剩余 C-H 键的 s 成分增加 \箭 键角变大(s 成分越多，键角越大，如 sp 的 180° \textgreater{} sp\textsuperscript{2} 的 120° \textgreater{} sp\textsuperscript{3} 的 109.5°)。
  \item
    选 \textbf{(a)} 变大。二氟甲烷中有两个 F，C-H 键的 s 成分进一步增加 \箭 H-C-H 键角继续增大。
  \end{enumerate}
\end{enumerate}

\begin{center}\rule{0.5\linewidth}{0.5pt}\end{center}

\subsection{八、延伸阅读}\label{ux516bux5ef6ux4f38ux9605ux8bfb}

\begin{itemize}
\tightlist
\item
  宋天佑《无机化学》§5 章(原子结构与周期律)
\item
  上海中学竞赛课程·化学·第一分册·第三讲
\end{itemize}

\begin{center}\rule{0.5\linewidth}{0.5pt}\end{center}

\end{document}
